\bookchapter{第六十三章：ビュー、状態、ナビゲーション}

\chapterintro{1ページだけのおもちゃなら、もう作れます。この章で解決するのは「成長」にまつわる2つの問題：複数のページの間をどう行き来するか（ナビゲーション）と、データがページの間をどう流れるか（状態管理）です。この章を終えれば、「リスト 詳細 編集」という定番のマルチページアプリの骨組みが組めるようになります。}

\booksection{状態管理の「データフロー全体図」}

SwiftUIの状態管理ツールは@Stateだけではありません。まずこの全体図を見て、「いつどれを使うか」をつかみましょう：

\begin{center}
\begin{tabularx}{\linewidth}{|X|X|X|X|}
\hline
\textbf{ツール} & \textbf{担当するもの} & \textbf{使う側} & \textbf{典型的な場面} \tabularnewline
\hline
@State & そのページだけの一時的な状態 & いまのビュー & カウンター、入力欄の文字 \tabularnewline
@Binding & ほかの状態を借りる & 子ビューから親ビューを書き換える & 自作コンポーネント \tabularnewline
@StateObject / @Observable & ページをまたいで共有するロジック & ビューモデル（ViewModel） & データ読み込み、ビジネスロジック \tabularnewline
@EnvironmentObject / .environment & アプリ全体で共有する状態 & アプリ全体 & ログイン状態、テーマ設定 \tabularnewline
\hline
\end{tabularx}
\end{center}

\begin{center}\small\color{gray}SwiftUIの状態管理ツール全体図（2026年、@Observableが主流）\end{center}

\begin{lstlisting}[language=Swift]
// iOS 17+のモダンな書き方：@Observableマクロ（新規プロジェクトにおすすめ）
import Observation

@Observable
class CounterModel {                 // ビューモデル
    var count = 0
    func increment() { count += 1 }
}

struct CounterView: View {
    @State private var model = CounterModel()   // モデルを保持

    var body: some View {
        VStack {
            Text("\(model.count)回押しました")
            Button("押して") { model.increment() }
        }
    }
}
\end{lstlisting}

初心者へのぶっきらぼうなアドバイス：ページの中だけで使うものは@State、ページをまたぐもの／ロジックを含むものは@Observableのクラスへ。古いチュートリアルの@StateObject/@ObservableObjectは前世代の書き方なので、読めればOK。新しいコードは上のルールに従いましょう。

\booksection{NavigationStack：ページ移動の「スタック」}

iOS 16以降、複数ページのナビゲーションにはNavigationStackを使います——お皿を積むようなもので、進むときに1枚ずつ積んでいき、戻るときは1枚ずつ取り出します：

\begin{lstlisting}[language=Swift]
struct HomeView: View {
    var body: some View {
        NavigationStack {
            List {
                NavigationLink("たろうの家計簿") {
                    DetailView(name: "たろうの家計簿")   // 行き先のページ
                }
                NavigationLink("はなこの家計簿") {
                    DetailView(name: "はなこの家計簿")
                }
            }
            .navigationTitle("家計簿リスト")   // 上の大きな見出し
        }
    }
}

struct DetailView: View {
    let name: String
    var body: some View {
        Text("\(name) の詳細ページ")
            .navigationTitle(name)
            .navigationBarTitleDisplayMode(.inline)
    }
}
\end{lstlisting}

NavigationLinkは「タップできる移動リンク」です。navigationTitleはナビゲーションバーのタイトルと戻るボタンを自動で生成します。プログラムで移動させたいとき（たとえば「保存」を押したら自動で戻る）は、@Environment(\textbackslash{}.dismiss)で現在のページを閉じます——この2つの合わせ技で、ナビゲーションのニーズの90％はカバーできてしまいます。

\booksection{ポップアップと下部タブ：sheetとTabView}

ページを「押し込む」以外にも、定番の操作が2つあります：

\begin{lstlisting}[language=Swift]
// 1. 下から半画面/全画面で表示（入力フォームの定番）
struct HomeView: View {
    @State private var showAdd = false
    var body: some View {
        Button("記帳する") { showAdd = true }
            .sheet(isPresented: $showAdd) {
                AddRecordView()     // 表示される入力ページ
            }
    }
}

// 2. 下部タブバー（アプリの「骨格」）
struct MainTabView: View {
    var body: some View {
        TabView {
            RecordListView()
                .tabItem { Label("家計簿", systemImage: "book") }
            StatsView()
                .tabItem { Label("統計", systemImage: "chart.pie") }
            SettingsView()
                .tabItem { Label("設定", systemImage: "gear") }
        }
    }
}
\end{lstlisting}

\begin{tcbinfo}{ワンポイント：どの場面でどれを使う？}
経験則：メインの流れはNavigationStackでページをpush（リスト詳細）、一時的な操作はsheetのポップアップ（追加·編集·確認）、並び立つ大ブロックはTabView（下部タブ）。この3つの組み合わせこそ、ほとんどのアプリの骨組みです。

\end{tcbinfo}

\booksection{List：iOSの「国民的リスト」}

iOSでいちばんよく使い、いちばん「Appleらしい」コンポーネントがListです——区切り線、スワイプ削除、タップ時のハイライトが、はじめから付いてきます：

\begin{lstlisting}[language=Swift]
struct RecordListView: View {
    @State private var records = ["昼ごはん 25.5", "地下鉄 6.0", "給料 +8000"]

    var body: some View {
        List {
            ForEach(records, id: \.self) { r in
                Text(r)
            }
            .onDelete { indexSet in          // 左スワイプで削除
                records.remove(atOffsets: indexSet)
            }
        }
    }
}
\end{lstlisting}

リスト＋NavigationLink＋sheet＋TabView——この4つを組み合わせれば、「帳簿リスト タップで詳細 ＋で入力ポップアップ」という完全なアプリの骨組みがもう組めます。次章の家計簿プロジェクトは、まさにこの組み方で作ります。

\begin{tcbpractice}{やってみよう：「3ページの骨組み」を組む}
\booksection{アプリのライフサイクル：フォアグラウンドとバックグラウンド}

AndroidのActivityのライフサイクルと同じように、iOSアプリにも「誕生、フォアグラウンド、バックグラウンド、終了」という状態の変化があります。SwiftUIではscenePhaseという環境値で観察します：

\begin{lstlisting}[language=Swift]
@main
struct MyApp: App {
    @Environment(\.scenePhase) private var phase

    var body: some Scene {
        WindowGroup {
            ContentView()
        }
        .onChange(of: phase) { _, newPhase in
            switch newPhase {
            case .active:    print("前面に戻る。更新を再開")
            case .inactive:  print("移行中（着信/バックグラウンド切替の瞬間）")
            case .background: print("バックグラウンドへ。下書きを保存")
            @unknown default: break
            }
        }
    }
}
\end{lstlisting}

典型的な使い方：フォアグラウンドに戻ったらデータを更新、バックグラウンドに入ったら下書きを保存。iOSはバックグラウンド動作にとても厳しく——普通のアプリはバックグラウンドに入るとすぐ停止（凍結）されるので、裏でこっそり重い処理をさせるのは期待しないでください。どうしても定期実行したいなら、「バックグラウンド常駐」ではなく第六十八章のローカル通知を使います。

ダミーデータ版の家計簿の骨組みを作りましょう：TabViewの3つのタブ（帳簿／統計／設定） 帳簿ページはListでダミーデータを表示し、行をタップしたら詳細へ（NavigationStack） 右上の「＋」ボタンで入力用のsheetを出す。データ保存も通信もなし、純粋に骨組みの練習——この骨組みができていれば、第六十五章で本物のデータを流し込むのは水到渠成です。

\end{tcbpractice}

\begin{tcbsummary}{本章のまとめ}
\begin{itemize}
  \item 状態管理の4ツール：@State（ページ内）、@Binding（借用）、@Observable（ロジック）、.environment（全体）
  \item 新規プロジェクトのビューモデルは@Observableマクロで書く。古いチュートリアルの@ObservableObjectは読めればOK
  \item ページ移動はNavigationStack。@Environment(\textbackslash{}.dismiss)でプログラムから戻る
  \item メインの流れはページをpush、一時的な操作はsheet、並び立つ大ブロックはTabView
  \item Listは「国民的リスト」：区切り線もスワイプ削除も、はじめから使える
\end{itemize}

\end{tcbsummary}
